\chapter{Introduction}
%%%%%%%%%%%%%%%%%%%%%%

Human-Computer Interaction (HCI) in high-performance computing and competitive esports places stringent demands on input peripherals. In environments where users execute several hundred Actions Per Minute (APM), actuation latency and threshold stability can affect both task performance and musculoskeletal load~\cite{esports_tendinopathy, wrist_fatigue_genre_2025}.

\section{The Evolution of Continuous Sensing Hardware}
Peripheral hardware has shifted from binary mechanical switches toward continuous sensing interfaces. In HCI, threshold selection and activation logic have long been recognized as key determinants of usability and speed \cite{touchpad_comparison, impact2018}. Recent Hall-effect and inductive devices expose depth-dependent actuation and reset behavior to firmware, including Rapid Trigger mechanisms for dynamic re-actuation \cite{wooting_rapid_trigger, pcgamer_hall_effect, drone_rapid_trigger}. Commercial systems such as HITS further combine inductive sensing with configurable haptic feedback \cite{logitech_superstrike_2026, logitech_hits_patent, fox_superstrike_2026, techradar_superstrike_2026}. Collectively, these developments increase control resolution but also increase dependence on well-calibrated software thresholds \cite{dynamic_keyboard, magnetic_tuning, attackshark_taphold}.

\section{The Limitations of Static Actuation Thresholds}
Despite these hardware leaps, the software interpreting this data remains limited. Current analog peripherals rely on static firmware thresholds that are agnostic to human biomechanical variance. Players exhibit unique control signatures, including micro-hesitations and varying velocity curves, which often clash with rigid hardware logic~\cite{impact2018}. When static thresholds are set too aggressively, they fail to filter out physiological noise; when set too conservatively, they introduce unnecessary mechanical pre-travel. This friction results in sub-optimal latency or a high frequency of unintended actuations during high-stress scenarios~\cite{touchpad_comparison, dynamic_keyboard}.

\section{Neuromotor Fatigue and Performance Decay}
These hardware trends also intersect with user health. Higher APM in esports has been associated with repetitive strain injuries and tendinopathies \cite{esports_tendinopathy, difrancisco2019managing, 1hp_esports_ergonomics}. Biomechanical studies report distinct fatigue patterns for keyboard and mouse interaction \cite{fatigue_flexor}, and repetitive clicking can degrade aiming accuracy while increasing forearm muscle load \cite{clicking_fatigue_ergonomics_2024}. High-APM genres are linked to elevated kinematic demand on wrist extensors \cite{wrist_fatigue_genre_2025}, with measurable neuromotor adaptation during fatiguing tasks \cite{heimhofer2024finger}. These findings motivate adaptive calibration that maintains responsiveness while reducing unintended actuation under fatigue.


\section{Toward Adaptive Intent Detection}
One longer-term direction is intent inference before full mechanical actuation. Prior work on surface electromyography (sEMG) and embedded TinyML suggests this is technically plausible~\cite{reactemg, lda_emg, tinyml_mcu_inference, singh2023edge}. However, this thesis focuses on a narrower and directly evaluable scope: offline calibration from depth telemetry already produced by the peripheral. On-device inference is treated as future work (Section~\ref{sec:conclusion-fw}). Productization and UX integration (including the UX Exploration v2 flow and eventual GHub/GHabits integration) are likewise positioned as future work.

\section{Thesis and Contributions}
The central hypothesis of this thesis is that per-user variation in click telemetry is structured enough to support automatic labeling, statistical modeling, and parameter optimization beyond static default settings. The contributions are:
\begin{itemize}
        \item A hybrid labeling pipeline for analog HID click intent that derives
            initial frame-level \texttt{MAKE\_ONSET}, \texttt{CLICK\_HOLD}, and \texttt{BREAK\_ONSET}
            annotations from algorithmic criteria (geometry, speed, and acceleration)
            and then improves boundary quality through model-assisted iteration.
    \item A four-class temporal sequence classifier with two model paths: a shallow
          path (Random Forest, Gradient Boosted Trees) operating on hand-crafted
          sliding-window features for sub-second per-session training, and a deep path
          (dilated residual FCN) operating on raw telemetry with a
          3{,}060-frame ($\approx$3\,s) learned receptive field.
    \item A firmware-faithful threshold optimizer that exhaustively evaluates the
          discrete (AP, RT) grid via a Rapid Trigger FSM simulator, returning the
          minimum-loss configuration and a Pareto frontier over detection rate and
          timing accuracy, replacing manual threshold tuning with a data-driven
          trade-off curve.
\end{itemize}

While not a primary algorithmic contribution, the work also includes a small HCI-facing component: purpose-built minigames that recreate common competitive-gaming input scenarios, providing a structured environment for users to generate calibration telemetry and receive a hardware recommendation.
The recommendation outputs are structured for interpretability (trade-off views and
comparative candidate points), so users can inspect why a suggested AP/RT point is proposed.

%%%%%%%%%%%%%%%%%%%%%%
